
\documentclass{article}
\usepackage{colortbl}
\usepackage{makecell}
\usepackage{multirow}
\usepackage{supertabular}

\begin{document}

\newcounter{utterance}

\centering \large Interaction Transcript for game `dond', experiment `semi\_en', episode 7 with reflexion{-}Qwen3.5{-}4B.
\vspace{24pt}

{ \footnotesize  \setcounter{utterance}{1}
\setlength{\tabcolsep}{0pt}
\begin{supertabular}{c@{$\;$}|p{.15\linewidth}@{}p{.15\linewidth}p{.15\linewidth}p{.15\linewidth}p{.15\linewidth}p{.15\linewidth}}
    \# & $\;$A & \multicolumn{4}{c}{Game Master} & $\;\:$B\\
    \hline

    \theutterance \stepcounter{utterance}  
    & & \multicolumn{4}{p{0.6\linewidth}}{
        \cellcolor[rgb]{0.9,0.9,0.9}{
            \makecell[{{p{\linewidth}}}]{
                \texttt{\tiny{[P1$\langle$GM]}}
                \texttt{You have attempted this same task before and did not succeed. These are the plans you wrote after those attempts.} \\
\texttt{Trial \#1: New plan: The previous attempt was aborted because I included explanatory text and reasoning before the actual message to the other player, which likely violated the rule "Do not include any response to the rules or text announcing your message." Additionally, I failed to strictly adhere to the format for the final agreement or proposal. The correct strategy is to send a single, concise message directly to the other player proposing the specific division of items without any preamble, analysis, or meta{-}commentary. I will propose taking the plate and one hammer, leaving the other hammer and two rings for the opponent, ensuring the message contains only the proposal statement.} \\
\\ 
\texttt{Trial \#2: New plan: The previous attempt was aborted because I sent a proposal message ("Proposal: 1 plate, 1 hammer") instead of a direct communication to the other player. The rules explicitly state that messages to the other player must not include the "Proposal:" prefix or square brackets, as those formats trigger the secret proposal mechanism. Since the game has already reached the time limit and the previous direct messages were likely ignored or misinterpreted due to the format, I must now submit a valid secret proposal using the exact required format: "{[}Proposal: <number> <object name>, <number> <object name>, ...{]}". I will calculate the optimal split based on my values (plate: 4, hammer: 2, ring: 1) to maximize my score. Given the opponent's likely strategy (they value rings at 1, so they might take them), I should propose taking the high{-}value items. A winning strategy is to propose taking the plate and both hammers (value 6) and leaving the two rings (value 2) for the opponent, or if the opponent takes rings, propose taking the plate and one hammer. I will format the final message strictly as "{[}Proposal: 1 plate, 2 hammers{]}" to ensure it is recognized as a secret proposal and not a direct message.} \\
\\ 
\texttt{Trial \#3: New plan: The previous attempt failed because I sent a direct message to the opponent ("I will take the plate and one hammer...") instead of submitting a secret proposal when the game state indicated it was time to finalize the agreement. The rules state that if proposals are complementary, both get points; otherwise, both get zero. Since the opponent already submitted a secret proposal, I must now submit a matching secret proposal to lock in the deal. I will analyze my value function (plate: 4, hammer: 2, ring: 1) and the opponent's likely move. If the opponent took the plate (my highest value), I should take the hammers. If they took the hammers, I should take the plate. To maximize my score, I will propose taking the plate and one hammer (total value 6) and implicitly leave the remaining items (one hammer, two rings) for the opponent, assuming the opponent's proposal complements this. I will format the message strictly as "{[}Proposal: 1 plate, 1 hammer{]}" to trigger the secret proposal mechanism and finalize the agreement.} \\
\\ 
\\ 
\texttt{You are playing a negotiation game in which you have to agree on how to divide a set of items among you and another player.} \\
\\ 
\texttt{Rules:} \\
\texttt{(a) You and the other player are given a set of items. Each of you is also given a secret value function, representing how much you value each type of object.} \\
\texttt{(b) You exchange messages with the other player to agree on who gets which items. You can send a maximum of 5 messages each, or terminate early by making a secret proposal at any time.} \\
\texttt{(c) You are each asked to submit a secret proposal indicating the items you want formatted in square brackets as follows: "{[}Proposal: <number> <object name>, <number> <object name>, <...>{]}"} \\
\texttt{(d) If your proposals are complementary, i.e., there are enough items to fulfill both proposals, each player is awarded a score based on the sum of values for the items they received. Otherwise, both of you get zero points.} \\
\texttt{(e) Your goal is to maximize the score you receive.} \\
\\ 
\texttt{Let us start.} \\
\\ 
\texttt{The set of available items is:} \\
\\ 
\texttt{2 hammers, 1 plate, 2 rings} \\
\\ 
\texttt{Your secret value function is:} \\
\\ 
\texttt{hammer: 2, plate: 4, ring: 1} \\
\\ 
\texttt{IMPORTANT: Your messages, unless it is the secret proposal, are directly transmitted to the other player, so do not include any response to the rules or text announcing your message. To make a secret proposal, use the indicated format. Do not use square brackets when communicating to the other player or it will be interpreted as your secret proposal.} \\
            }
        }
    }
    & \\ \\

    \theutterance \stepcounter{utterance}  
    & \multicolumn{4}{p{0.6\linewidth}}{
        \cellcolor[rgb]{0.9,0.9,0.9}{
            \makecell[{{p{\linewidth}}}]{
                \texttt{\tiny{[P1$\rangle$GM]}}
                \texttt{I will take the plate and one hammer.} \\
            }
        }
    }
    & & \\ \\

    \theutterance \stepcounter{utterance}  
    & & & \multicolumn{2}{p{0.3\linewidth}}{
        \cellcolor[rgb]{0.9,0.9,0.9}{
            \makecell[{{p{\linewidth}}}]{
                \texttt{\tiny{[GM$|$GM]}}
                \texttt{continue} \\
            }
        }
    }
    & & \\ \\

    \theutterance \stepcounter{utterance}  
    & & & \multicolumn{2}{p{0.3\linewidth}}{
        \cellcolor[rgb]{0.9,0.9,0.9}{
            \makecell[{{p{\linewidth}}}]{
                \texttt{\tiny{[GM$|$GM]}}
                \texttt{I will take the plate and one hammer.} \\
            }
        }
    }
    & & \\ \\

    \theutterance \stepcounter{utterance}  
    & & \multicolumn{4}{p{0.6\linewidth}}{
        \cellcolor[rgb]{0.9,0.9,0.9}{
            \makecell[{{p{\linewidth}}}]{
                \texttt{\tiny{[GM$\rangle$P2]}}
                \texttt{You have attempted this same task before and did not succeed. These are the plans you wrote after those attempts.} \\
\texttt{Trial \#1: New plan: The previous attempt failed because I sent "Proposal: 2 rings" as a direct message to the other player instead of using the required secret proposal format. The rules state that direct messages should not include square brackets, but the secret proposal itself must be formatted exactly as "{[}Proposal: <number> <object name>, ...{]}". I misunderstood that the square brackets were only forbidden in *communication* (direct messages), not in the *proposal* format itself. Next time, I will ensure that any message containing the word "Proposal" is strictly enclosed in square brackets (e.g., "{[}Proposal: 2 rings{]}") to trigger the secret proposal mechanism, and I will avoid sending unformatted text that looks like a proposal. Additionally, I will verify that the items listed in the proposal actually exist in the available set (2 rings are available) before submitting.} \\
\\ 
\texttt{Trial \#2: New plan: The previous attempt failed because I made a direct offer ("I will take the plate and one hammer...") which was interpreted as a secret proposal due to the lack of square brackets, but the content was incomplete or not complementary to a valid counter{-}proposal in the game's logic flow. I need to start with a direct negotiation message to establish a trade, rather than immediately jumping to a secret proposal. My strategy will be to propose a specific exchange (e.g., "I give you the plate, you give me the hammer") in plain text. Once the other player responds, I will analyze their offer. If their offer is complementary to mine (i.e., we can split the items such that both get what they value), I will then submit a secret proposal containing the agreed{-}upon split. If their offer is not complementary, I will adjust my direct message to propose a different split that maximizes my value (plate: 2, ring: 4) while ensuring the total items match the available set. I will strictly avoid sending square{-}bracketed proposals until I have confirmed a complementary agreement through direct communication.} \\
\\ 
\\ 
\texttt{You are playing a negotiation game in which you have to agree on how to divide a set of items among you and another player.} \\
\\ 
\texttt{Rules:} \\
\texttt{(a) You and the other player are given a set of items. Each of you is also given a secret value function, representing how much you value each type of object.} \\
\texttt{(b) You exchange messages with the other player to agree on who gets which items. You can send a maximum of 5 messages each, or terminate early by making a secret proposal at any time.} \\
\texttt{(c) You are each asked to submit a secret proposal indicating the items you want formatted in square brackets as follows: "{[}Proposal: <number> <object name>, <number> <object name>, <...>{]}"} \\
\texttt{(d) If your proposals are complementary, i.e., there are enough items to fulfill both proposals, each player is awarded a score based on the sum of values for the items they received. Otherwise, both of you get zero points.} \\
\texttt{(e) Your goal is to maximize the score you receive.} \\
\\ 
\texttt{Let us start.} \\
\\ 
\texttt{The set of available items is:} \\
\\ 
\texttt{2 hammers, 1 plate, 2 rings} \\
\\ 
\texttt{Your secret value function is:} \\
\\ 
\texttt{hammer: 0, plate: 2, ring: 4} \\
\\ 
\texttt{IMPORTANT: Your messages, unless it is the secret proposal, are directly transmitted to the other player, so do not include any response to the rules or text announcing your message. To make a secret proposal, use the indicated format. Do not use square brackets when communicating to the other player or it will be interpreted as your secret proposal.} \\
\\ 
\\ 
\texttt{I will take the plate and one hammer.} \\
            }
        }
    }
    & \\ \\

    \theutterance \stepcounter{utterance}  
    & & & \multicolumn{4}{p{0.6\linewidth}}{
        \cellcolor[rgb]{0.9,0.9,0.9}{
            \makecell[{{p{\linewidth}}}]{
                \texttt{\tiny{[GM$\langle$P2]}}
                \texttt{I will take the plate and one hammer.} \\
            }
        }
    }
     \\ \\

    \theutterance \stepcounter{utterance}  
    & & & \multicolumn{2}{p{0.3\linewidth}}{
        \cellcolor[rgb]{0.9,0.9,0.9}{
            \makecell[{{p{\linewidth}}}]{
                \texttt{\tiny{[GM$|$GM]}}
                \texttt{continue} \\
            }
        }
    }
    & & \\ \\

    \theutterance \stepcounter{utterance}  
    & & & \multicolumn{2}{p{0.3\linewidth}}{
        \cellcolor[rgb]{0.9,0.9,0.9}{
            \makecell[{{p{\linewidth}}}]{
                \texttt{\tiny{[GM$|$GM]}}
                \texttt{I will take the plate and one hammer.} \\
            }
        }
    }
    & & \\ \\

    \theutterance \stepcounter{utterance}  
    & & \multicolumn{4}{p{0.6\linewidth}}{
        \cellcolor[rgb]{0.9,0.9,0.9}{
            \makecell[{{p{\linewidth}}}]{
                \texttt{\tiny{[P1$\langle$GM]}}
                \texttt{I will take the plate and one hammer.} \\
            }
        }
    }
    & \\ \\

    \theutterance \stepcounter{utterance}  
    & \multicolumn{4}{p{0.6\linewidth}}{
        \cellcolor[rgb]{0.9,0.9,0.9}{
            \makecell[{{p{\linewidth}}}]{
                \texttt{\tiny{[P1$\rangle$GM]}}
                \texttt{{[}Proposal: 1 plate, 1 hammer{]}} \\
            }
        }
    }
    & & \\ \\

    \theutterance \stepcounter{utterance}  
    & & & \multicolumn{2}{p{0.3\linewidth}}{
        \cellcolor[rgb]{0.9,0.9,0.9}{
            \makecell[{{p{\linewidth}}}]{
                \texttt{\tiny{[GM$|$GM]}}
                \texttt{proposal} \\
            }
        }
    }
    & & \\ \\

    \theutterance \stepcounter{utterance}  
    & & & \multicolumn{2}{p{0.3\linewidth}}{
        \cellcolor[rgb]{0.9,0.9,0.9}{
            \makecell[{{p{\linewidth}}}]{
                \texttt{\tiny{[GM$|$GM]}}
                [1, 1, 0]
            }
        }
    }
    & & \\ \\

    \theutterance \stepcounter{utterance}  
    & & \multicolumn{4}{p{0.6\linewidth}}{
        \cellcolor[rgb]{0.9,0.9,0.9}{
            \makecell[{{p{\linewidth}}}]{
                \texttt{\tiny{[GM$\rangle$P2]}}
                \texttt{The other player has submitted a secret proposal. It is now your turn to submit a proposal of your own.} \\
            }
        }
    }
    & \\ \\

    \theutterance \stepcounter{utterance}  
    & & & \multicolumn{4}{p{0.6\linewidth}}{
        \cellcolor[rgb]{0.9,0.9,0.9}{
            \makecell[{{p{\linewidth}}}]{
                \texttt{\tiny{[GM$\langle$P2]}}
                \texttt{{[}Proposal: 1 ring, 1 hammer{]}} \\
            }
        }
    }
     \\ \\

    \theutterance \stepcounter{utterance}  
    & & & \multicolumn{2}{p{0.3\linewidth}}{
        \cellcolor[rgb]{0.9,0.9,0.9}{
            \makecell[{{p{\linewidth}}}]{
                \texttt{\tiny{[GM$|$GM]}}
                \texttt{proposal} \\
            }
        }
    }
    & & \\ \\

    \theutterance \stepcounter{utterance}  
    & & & \multicolumn{2}{p{0.3\linewidth}}{
        \cellcolor[rgb]{0.9,0.9,0.9}{
            \makecell[{{p{\linewidth}}}]{
                \texttt{\tiny{[GM$|$GM]}}
                [1, 0, 1]
            }
        }
    }
    & & \\ \\

    \theutterance \stepcounter{utterance}  
    & & & \multicolumn{2}{p{0.3\linewidth}}{
        \cellcolor[rgb]{0.9,0.9,0.9}{
            \makecell[{{p{\linewidth}}}]{
                \texttt{\tiny{[GM$|$GM]}}
                [[1, 1, 0], [1, 0, 1]]
            }
        }
    }
    & & \\ \\

\end{supertabular}
}

\end{document}
